%% kapa: exam=1 round=19 year=2008 period=2 form=A subject=economics subject_ko=경제원론 seq=35 number=35
%% source: Q-Net 감정평가사 기출문제 (cst00309) artlSeq=1001962 file=extracted/1008/1001962/2008년 제19회 감정평가사 1차시험 기출문제/2교시.hwp
%% conversion: kordoc 4.24.51 -> tools/gen1.py (auto); review=proofread-sonnet; flags=tikz
\begin{question}{35}
다음 그림과 같이 미숙련 노동시장의 균형임금은 시간당 5,000원으로 나타났다. 최저임금을 시간당 7,000원으로 설정하면, 1일의 경제적 순손실(deadweight loss)은 얼마인가?
\begin{center}
\begin{tikzpicture}[scale=0.55, font=\small]
  \draw (0,0) -- (0,10.8);
  \draw (0,0) -- (10.8,0);
  \node[above] at (1.2,10.9) {임금(원/시간)};
  \node[left] at (0,10) {10,000};
  \node[left] at (0,5) {5,000};
  \node[below left] at (0,0) {0};
  \node[below] at (5,-0.1) {5,000};
  \node[below] at (10,-0.1) {10,000};
  \node[right] at (10.9,0) {고용량(시간/일)};
  \draw (0,10) -- (10,0);
  \draw (0,0) -- (10,10);
  \draw[dotted] (0,5) -- (5,5) -- (5,0);
  \node[right] at (10.1,10.1) {노동공급};
  \node[right] at (8.3,2.1) {노동수요};
\end{tikzpicture}
\end{center}
\begin{choices}
  \item 0원
  \item 200만 원
  \item 400만 원
  \item 1,250만 원
  \item 2,500만 원
\end{choices}
\end{question}
\answer{3}
